\section{SPECIFIKÁCIÓ}\label{specifikuxe1ciuxf3}

\subsection{Program funkciói}\label{program-funkciuxf3i}
\subsubsection{\texorpdfstring{ Program célja}{ Program célja}}\label{program-cuxe9lja}

A projektem célja, hogy egy az autóra elejére, a szélvédő mögé és a
visszapillantó tükör közelében rögzített, monokuláris kamera
felviteléből, miközben az autó mozog az úton, a kamera származó
videófolyamot feldolgozva, a képkockákon lévő kátyúkat felismerni,
detektálni (képen lévő helyüket bekeretezni). A kátyú detektálását
követően azok pontosabb helyének meghatározása az autó orrához képest.
Ez alatt azt értem, hogy a kamerakép 2D koordináta rendszeréből a 3D
világ koordináta rendszerébe való konverziót. A lyukkamera modellre
építve, szeretném meghatározni a kátyú oldalirányú pozícióját az úton
belül, a felvételt készítő autóhoz képest. Mindezt valós időben, azaz
miközben a jármű előrefelé halad az úton a szoftvernek az autó előtti
területet kell vizsgálnia, és előre detektálni a kátyúkat és azok
oldalirányú pozícióit, mielőtt az autó oda érne a kátyús útszakasz
részhez. Tehát egyfajta szoftveres előre tekintés az úton, ami nem a
lámpákat, közlekedési táblákat és nem a felfestéseket figyeli
útburkolaton, hanem az úthibákat, azon belül is kifejezetten a kátyúkra
fókuszál.

Meglátásom szerint egy ilyen rendszerrel alapot lehetne teremteni
autonóm kátyú kikerülő manőverek végrehajtására egyfajta vezetéstámogató
rendszer keretein belül. Ez nemcsak a járművek állapotának megóvásában
jelentene előnyt, hanem növelné a közlekedés biztonságát és nem utolsó
sorban a menetkomfortot is. Ez indokolj azt a döntésemet, hogy csak egy
darab monokuláris kamra képpel dolgozok. Mert a jelenlegi modern autók
esetében, sok advanced driver assistance systems (ADAS) funkció egy
kamera kép feldolgozásával van megoldva.

Szeretném ki hangsúlyozni, hogy a projekt keretein belül elsősorban csak
a kátyúk magabiztos felismerése, azok detektálása és az úton belül
betöltött oldalirányú helyzetük minél pontosabb meghatározása áll, a
felvételt készítő járműhöz képest, valós időben. Tehát a rendszer nem
foglalkozik a járműdinamikai beavatkozással (kormányzás), kizárólag a
környezetérzékeléssel. Az esetleges vezetéstámogató rendszer rész egy
ráépülő önvezetéssel kapcsolatos továbbfejlesztési ötlet, ami a
motivációt adja a projekthez.

\subsubsection{Elérhető funkciók
részletesen}\label{eluxe9rhetux151-funkciuxf3k-ruxe9szletesen}

\paragraph{Videófolyam feldolgozás és
előfeldolgozás}\label{videuxf3folyam-feldolgozuxe1s-uxe9s-elux151feldolgozuxe1s}

A rendszer belépési pontja a valós idejű képfeldolgozó modul, amely a
szélvédő mögött elhelyezett monokuláris kamera jelét fogadja.

Első lépésben a szoftver folyamatosan olvassa és digitalizálja a kamera
által biztosított 2D videófolyamot. A feldolgozási lánc biztosítja, hogy
a rendszer képes legyen tartani a bemeneti forrás sebességét
(szinkronizált 30 FPS), ezáltal garantálva a folyamatos, akadozásmentes
működést nagy sebességű haladás esetén is.

Majd mielőtt a detektáló algoritmus lefutna, a rendszer előfeldolgozást
végez az egyes képkockákon, hogy a nappali, természetes fényviszonyok
mellett a lehető legkontrasztosabb képet biztosítsa a kátyúk
körvonalainak kiemeléséhez.

\paragraph{Kátyú detektálás, objektum
felismerés}\label{kuxe1tyuxfa-detektuxe1luxe1s-objektum-felismeruxe9s}

Ez a modul a rendszer szeme, amely a vizuális információk értelmezéséért
felelős, gépi tanulás segítségével.

A szoftver képes azonosítani az aszfaltozott útfelületen felszínen
található durva, egyenletlen tál alakú mélyedés, azaz kátyúkat. A
funkció kritikus eleme a megkülönböztetés képesség, tehát az algoritmus
robusztusan különíti el a kátyúkat az egyéb vizuális zajoktól. A
rendszer nem tekinti kátyúnak az árnyékokat, az útburkolati jeleket, az
aszfalt foltszerű javításait vagy a repedéseket, így minimalizálva a
false positive-ok számát.

A funkció előre tekintő detektálást végez, a jármű előtti teret
pásztázva, már 30 méteres távolságból azonosítsa a potenciális
veszélyforrásokat. Ez az előre tekintés biztosítja a szükséges
reakcióidőt a sofőr vagy egy későbbi autonóm rendszer számára.

\paragraph{\texorpdfstring{ Térbeli Lokalizáció és
Koordináta-transzformáció}{ Térbeli Lokalizáció és Koordináta-transzformáció}}\label{tuxe9rbeli-lokalizuxe1ciuxf3-uxe9s-koordinuxe1ta-transzformuxe1ciuxf3}

A rendszer legfontosabb matematikai modulja, amely a detektált
objektumokat elhelyezi a valós térben a járműhöz viszonyítva.

2D-ből 3D transzformációhoz a lyukkamera-modell (pinhole camera model)
matematikai összefüggéseit alkalmazva a szoftver a sík képi
koordinátákat átkonvertálja a 3D világkoordináta rendszerbe. Ez a
számítás teszi lehetővé, hogy a képernyőn látott foltból valós fizikai
távolságadat legyen.

Laterális (oldalirányú) pozíció meghatározásánál a rendszer kiszámítja a
detektált kátyú oldalirányú elhelyezkedését a jármű hossztengelyéhez
(orrához) képest. A funkció célja, hogy centiméteres pontossággal
megadja, az autó nyomtávjához viszonyítva hol helyezkedik el a kátyú az
úton.

\paragraph{Valós Idejű Vizualizáció (Ember-Gép
Interfész)}\label{valuxf3s-idejux171-vizualizuxe1ciuxf3-ember-guxe9p-interfuxe9sz}

A szoftver vizuális visszajelzést biztosít a feldolgozás eredményéről.
Kiterjesztett valóság jelleggel, a futtató eszköznek (laptop) a
kijelzőjén megjeleníti az élő videóképet grafikus határoló keretekkel
(bounding box) jelöli meg a detektált kátyúkat. A határoló dobozok
mellett a rendszer szöveges formában, valós időben megjeleníti a
\hyperref[tuxe9rbeli-lokalizuxe1ciuxf3-uxe9s-koordinuxe1ta-transzformuxe1ciuxf3]{3.1.2.3}
pontban számított távolság adatokat. Ez lehetővé teszi a rendszer
működésének azonnali ellenőrzését működés közben.

\subsection{Követelmények}\label{kuxf6vetelmuxe9nyek}

\subsubsection{Követelmények
kifejtve}\label{kuxf6vetelmuxe9nyek-kifejtve}

\paragraph{Funkcionális
követelmények}\label{funkcionuxe1lis-kuxf6vetelmuxe9nyek}

Ezen követelmények határozzák meg a rendszer alapvető működési
mechanizmusait, a bemeneti adatok kezelésétől kezdve a feldolgozáson át
egészen az eredmények megjelenítéséig.

Bemeneti videófolyam feldolgozása A rendszer elsődleges bemeneti
interfészének képesnek kell lennie a jármű szélvédője mögé, a
visszapillantó tükör környezetében rögzített monokuláris kamera jelének
fogadására. A szoftvernek folyamatosan, megszakítás nélkül kell
feldolgoznia a kamera által biztosított, kétdimenziós videófolyamot. Ez
a funkció képezi a környezetérzékelés alapját, így elengedhetetlen, hogy
a rendszer a képkockák egymásutániságát időben helyesen kezelje,
biztosítva ezzel a későbbi modulok számára szükséges adatfolyamot.

Szelektív detektálás és objektumosztályozás A képfeldolgozó
algoritmusnak magas megbízhatósággal kell azonosítania az útfelületen
található geometriai mélyedéseket, azaz a kátyúkat. A funkció kritikus
eleme a szelektivitás, a rendszernek képesnek kell lennie
megkülönböztetni a valódi veszélyt jelentő úthibákat az egyéb vizuális
zajoktól. A detektálási logikának ki kell szűrnie a false positive-okat,
amelyeket az útburkolati jelek, az aszfalt repedései, a foltszerű
javítások, vagy a környezeti tárgyak (fák, épületek) által vetett
árnyékok okozhatnak. Kizárólag a jármű futóművére veszélyes mélyedések
kerülhetnek detektálásra.

Valós idejű vizualizáció és visszajelzés, a fejlesztés, tesztelés és a
későbbi felhasználói felügyelet támogatása érdekében a rendszernek
vizuális visszajelzést kell biztosítania a feldolgozás eredményéről. A
feldolgozást végző eszköz (laptop) képernyőén megjelenő élő videóképen a
detektált kátyúkat egyértelműen, a kiterjedésüket jelölő határoló
dobozzal (bounding box) kell kiemelni. Ezen felül a vizualizációnak
tartalmaznia kell a számított numerikus adatokat is, különös tekintettel
a kátyú járműtől mért távolságára és az oldalirányú pozíciójára, ezzel
segítve a rendszer pontosságának szubjektív és objektív ellenőrzését.

Működési tartomány a rendszer alkalmazhatósága a projekt jelenlegi
fázisában meghatározott útviszonyokra és járműmozgásra korlátozódik. A
detektálási algoritmusnak aszfaltozott burkolattal ellátott utakon kell
üzembiztosan működnie. A geometria számítások egyszerűsítése és a
lyukkamera-modell érvényessége érdekében a rendszernek elsősorban
egyenes vonalú haladáskor kell helytállnia. A specifikáció kizárja a
rendszer működéséből a tolatási manővereket, valamint a hirtelen
irányváltással járó éles kanyarokat, mivel ezekben az esetekben a
kamerakép és a járműmozgás vektora közötti korreláció jelentősen eltér a
modell feltételezéseitől.

\paragraph{Minőségi
követelmények}\label{minux151suxe9gi-kuxf6vetelmuxe9nyek}

A minőségi követelmények definiálják azokat a teljesítménymutatókat és
korlátokat, amelyeken belül a funkcionális követelményeknek teljesülniük
kell a biztonságos és hatékony működés érdekében.

Valós idejű működés, rendszer reakcióképességének kulcsa, hogy a teljes
feldolgozási lánc, beleértve a kép előfeldolgozását, a detektálást, és a
matematikai lokalizációt, képes legyen követni a bemeneti videó
sebességét. Mivel a kamera 30 képkocka/másodperc (FPS) sebességgel
rögzít, a szoftvernek rendelkeznie kell azzal a számítási kapacitással,
hogy egyetlen képkocka feldolgozása ne haladja meg a \textasciitilde33
milliszekundumot. Ez biztosítja, hogy ne alakuljon ki késleltetés, ami a
gyorsan mozgó jármű esetén adatvesztéshez és(vagy) elkésett
detektáláshoz vezetne.

Detektálási távolság és előre tekintés, a szoftvernek képesnek kell
lennie arra, hogy a kátyúkat 30 méteres távolságból azonosítsa. Ez a
távolság elegendő időt biztosítana egy elméleti autonóm rendszernek vagy
a sofőrnek a kikerülő manőver megtervezésére és végrehajtására, még
nagyobb sebesség mellett is.

Támogatott sebességtartomány A rendszernek le kell fednie a városi és az
országúti közlekedés jellemző sebességtartományait. A képfeldolgozó
algoritmusoknak robusztusnak kell lenniük a mozgási elmosódással
szemben, biztosítva a stabil detektálást álló helyzettől (0 km/h)
egészen 70 km/h sebességig. A szoftver működése nem degradálódhat a
jármű sebességének növekedésével ezen a tartományon belül.

Lokalizációs pontosság a projekt egyik legfontosabb műszaki kihívása a
2D képpontok 3D térbeli koordinátákká való konvertálása. A járműhöz
közel, 10 méteres távolságon belül a kátyú oldalirányú pozíciójának
becslési hibája nem haladhatja meg a ±15 cm-t. Ez a pontosság kritikus
fontosságú, mivel egy átlagos gumiabroncs szélessége is összemérhető
ezzel az értékkel.

A minél jobb környezeti robusztussághoz rengeteg adatra és számítási
kapacitásra van szükség. Ezért a szoftvernek csak nappali, természetes
fényviszonyok mellett kell megbízhatóan üzemelnie. Ez magában foglalja a
napsütéses, felhős vagy változóan árnyékos időszakokat is, azonban a
jelenlegi fázisban nem terjed ki az éjszakai vezetésre vagy az extrém
időjárási körülmények (köd, szakadó eső, havas táj) kezelésére.

\subsubsection{\texorpdfstring{Követelmények felsorolás jelleggel
}{Követelmények felsorolás jelleggel }}\label{kuxf6vetelmuxe9nyek-felsoroluxe1s-jelleggel}

\paragraph{Funkcionális
követelmények}\label{funkcionuxe1lis-kuxf6vetelmuxe9nyek-1}

\begin{enumerate}
\def\labelenumi{\arabic{enumi}.}
\item
  Bemenet: A rendszernek fel kell dolgoznia a szabványos, monokuláris
  kamera által biztosított 2D videófolyamot.
\item
  Szelektív detektálás: A szoftvernek képesnek kell lennie a kátyúk
  megkülönböztetésére egyéb útburkolati jelektől (felfestések,
  repedések, foltok, árnyékok).
\item
  Vizualizáció: A detektált hibákat a képernyőn határoló dobozzal és a
  számított távolságadatok megjelenítésével kell jelölni.
\item
  Működési tartomány: A detektálásnak aszfaltozott egyenes úton, egyenes
  vonalú haladás esetén kell működnie, tolatás és éles kanyarok esetén
  nem.
\end{enumerate}

\paragraph{Minőségi
követelmények}\label{minux151suxe9gi-kuxf6vetelmuxe9nyek-1}

\begin{enumerate}
\def\labelenumi{\arabic{enumi}.}
\item
  Valós idejűség: A teljes feldolgozási láncnak (pre-processing,
  detektálás, lokalizáció) követnie kell a bemeneti 30FPS videó
  sebességét.
\item
  Detektálási távolság: A rendszernek legalább 30 méter távolságból fel
  kell ismernie a kátyúkat.
\item
  Sebességtartomány: A funkciónak 0--70 km/h járműsebesség mellett
  stabilan kell működnie.
\item
  Lokalizációs pontosság: A kátyú oldalirányú pozíciójának becslési
  hibája a járműtől 10 méteren belül nem haladhatja meg a ±15 cm-t.
\item
  Robusztusság: A rendszernek nappali, természetes fényviszonyok mellett
  kell üzembiztosan működnie.
\end{enumerate}

\subsection{\texorpdfstring{Megvalósítási módszertan
}{Megvalósítási módszertan }}\label{megvaluxf3suxedtuxe1si-muxf3dszertan}

A rendszer komplexitása indokolta a hibrid fejlesztési módszertan
alkalmazását, amely ötvözi a modern, adat vezérelt mesterséges
intelligenciát a klasszikus, determinisztikus matematikai modellezéssel.
A fejlesztés során két eltérő, de egymást kiegészítő stratégiát
használok

\subsubsection{\texorpdfstring{Környezetérzékelés: Iteratív Kísérleti
Fejlesztés
}{Környezetérzékelés: Iteratív Kísérleti Fejlesztés }}\label{kuxf6rnyezetuxe9rzuxe9keluxe9s-iteratuxedv-kuxedsuxe9rleti-fejlesztuxe9s}

A kátyúk vizuális felismeréséért felelős modul létrehozása iteratív
(ciklikus) fejlesztési modell alapján történt, amely a gépi tanuláson
alapuló projektekre jellemző. A folyamat az adatgyűjtés, annotálás,
modelltanítás és kiértékelés (validáció) ismétlődő köreiből állt. A
teszteredmények folyamatos visszacsatolása alapján finomhangolm a modell
hiperparamétereit.

\subsubsection{Lokalizáció: Analitikus Modellalapú
Rekonstrukció}\label{lokalizuxe1ciuxf3-analitikus-modellalapuxfa-rekonstrukciuxf3}

A detektált hibák laterális pozíciójának és távolságának
meghatározásához deduktív, matematikai modellalapú módszertant
alkalmazok. Ez a megközelítés nem tanuláson, hanem szigorú geometriai
szabályokon, konkrétan a lyukkamera modell alkalmazásán alapul. A
módszer a síkút feltételezésre építve. A számítások pontosságát a kamera
belső és külső paramétereinek előzetes kalibrációja biztosítja.